\subsection{Sorting scalar contacts with a total oscillation bound}

We first give a deterministic statement.  In particular, its constant does
not depend on the degrees of the contact graph.  A \emph{minimum chord} of
a continuous function $f$ is a pair $(u,v)$ such that
\[
 f(u)=f(v)=\min_{[u\wedge v,u\vee v]}f.
\]
All graphs in the next Lemma are finite.  Parallel edges and repeated chord
endpoints are allowed.

\begin{lemma}[Sorted scalar contacts]\label{arm:sorted-scalar-contacts}
Let $f$ be continuous on a compact interval containing pairwise disjoint
nonempty compact intervals $I_1,\ldots,I_m$, ordered from left to right.
Choose $t_v\in I_v$ and write
\[
 h_v=f(t_v),\qquad \omega_v=\max_{I_v}f-\min_{I_v}f.
\]
Let $G$ be a graph on these intervals, each of whose edges is witnessed by
a minimum chord with one endpoint in each incident interval.  Let $C$ be
a connected component of $G$, and list its vertices in increasing order
as $v_1,\ldots,v_k$.  For $1\le i<k$ set
\begin{equation}\label{arm:sorted-contact-errors}
 \mu_i=\min_{[t_{v_i},t_{v_{i+1}}]} f,
 \qquad E_i=h_{v_i}+h_{v_{i+1}}-2\mu_i.
\end{equation}
Then $E_i\ge0$,
\begin{equation}\label{arm:sorted-contact-walls}
 |h_{v_{i+1}}-h_{v_i}|\le E_i,
 \qquad
 \inf_{[t_{v_i},t_{v_{i+1}}]}f
 \ge \max\{h_{v_i},h_{v_{i+1}}\}-E_i,
\end{equation}
and
\begin{equation}\label{arm:sorted-contact-total}
 \sum_{i=1}^{k-1}E_i\le 6\sum_{v\in C}\omega_v.
\end{equation}
Thus the wall bound also holds on the gap between the two intervals, or
on any part of that gap obtained by deleting other intervals.  The same
bound can be summed over all components of $G$.
\end{lemma}

\begin{proof}
Discard any loops of $G$, which do not change its connected components.
We construct a finite metric tree directly from the minima of $f$.  This
also explains the contour order used in the proof.

Let $[A,B]$ contain all the intervals, and let
$m_*=\min_{[A,B]}f$.  For each $I_v=[a_v,b_v]$, choose a point $r_v$ at
which $f$ attains its minimum on $I_v$.  Let $\mathcal T$ be the finite set
consisting of $A,B$, all $a_v,b_v,r_v,t_v$, all chord endpoints, and a
point attaining $m_*$.  Coincident marked times occur only once in this
set.  Write
\[
 M(s,t)=\min_{[s\wedge t,s\vee t]}f\qquad(s,t\in\mathcal T).
\]
For each $t\in\mathcal T$ take a vertical segment
$\{t\}\times[m_*,f(t)]$.  Identify $(s,z)$ with $(t,z)$ exactly when
$z\le M(s,t)$.  This is an equivalence relation: the only nontrivial
check follows from
\[
 M(s,u)\ge\min\{M(s,t),M(t,u)\},
\]
because the interval between $s$ and $u$ is contained in the union of the
intervals between $s,t$ and between $t,u$.

The quotient, with distance measured along the vertical segments, is a
finite rooted metric tree $T$.  To verify this explicitly, adjoin the
marked times in increasing order.  The segment for a new time $t_j$
shares with the previously constructed tree precisely the initial
segment up to height $M(t_{j-1},t_j)$: indeed
\[
 \max_{i<j}M(t_i,t_j)=M(t_{j-1},t_j).
\]
Its remaining segment is therefore attached at a single point.  This
induction constructs a tree, allowing a new segment to have length zero.
Declare all marked tips and all attachment points to be vertices of $T$.
All its segments start at the common root of height $m_*$.  Denote the
top of the segment for $t$ by $\pi(t)$.  The common part of the two root
paths for $\pi(s),\pi(t)$ ends at height $M(s,t)$, so
\begin{equation}\label{arm:finite-contour-tree-distance}
 d_T(\pi(s),\pi(t))=f(s)+f(t)-2M(s,t).
\end{equation}
In particular, the endpoints of every minimum chord have the same image
in $T$.

For $I_v=[a_v,b_v]$, put $m_v=f(r_v)$ and define the subtree
\[
 L_v=[\pi(a_v),\pi(b_v)]_T
       \mathbin{\cup}[\pi(r_v),\pi(t_v)]_T.
\]
Here and below $[x,y]_T$ denotes the unique tree arc between $x$ and $y$.
The point $\pi(r_v)$ lies on the first arc: it is the common ancestor of
$\pi(a_v)$ and $\pi(b_v)$ at height $m_v$.  It is also an ancestor of
$\pi(t_v)$.  Thus $L_v$ is connected.  If $\ell$ denotes the length of a
finite metric tree, then
\begin{align}\label{arm:box-skeleton-length}
 \ell(L_v)
 &\le f(a_v)+f(b_v)-2m_v+f(t_v)-m_v\\
 &\le 3\omega_v.\notag
\end{align}
Overlapping arcs only reduce the first quantity.

Every chord contact in $I_v$ with a time outside $I_v$ belongs to
$[\pi(a_v),\pi(b_v)]_T$.  To see this, let $u\in I_v$ be the contact
and let $w<a_v$ be the other endpoint.  The chord condition implies
\[
 M(a_v,u)=f(u),\qquad m_v\le f(u)\le f(a_v).
\]
Hence $\pi(u)$ is the point of height $f(u)$ on the root path to
$\pi(a_v)$, between $\pi(r_v)$ and $\pi(a_v)$.  If instead $w>b_v$,
the same argument uses $b_v$.  This proves the claim, including contacts
at interval endpoints.

For each edge of $G$ joining $v$ to $w$, its common chord image therefore
belongs to $L_v\cap L_w$.  Graph connectivity shows that
\[
 U_C=\bigcup_{v\in C}L_v
\]
is connected.  A connected subset of a tree contains the unique arc
between any two of its points: removing an omitted point of that arc
would separate the two points.  Consequently $U_C$ is a subtree and
contains the finite hull $H$ of
$\pi(t_{v_1}),\ldots,\pi(t_{v_k})$.  Moreover,
\begin{equation}\label{arm:connected-box-skeleton-length}
 \ell(H)\le\ell(U_C)\le3\sum_{v\in C}\omega_v.
\end{equation}

We next justify the contour tour bound for the chronologically ordered
representatives.  Subdivide $H$ at every vertex of $T$ lying in $H$,
including all representative images; there are finitely many resulting
edges.  For a point $x$ in the interior of such an edge, let $\lambda$ be
its absolute height.  Among the components of $T\setminus\{x\}$, call the one
not containing the root the distal component.  The representative
indices whose images belong to this component form a consecutive block.
Indeed, if the images at times $s<u$ belong to it, then their common
ancestor has absolute height strictly greater than $\lambda$.
Thus $M(s,u)>\lambda$.  For any marked time $t\in[s,u]$,
\[
 M(s,t)\ge M(s,u)>\lambda,
\]
which puts $\pi(t)$ in the same distal component.  The strict inequality
is valid because $x$ is in an edge interior and all marked images are
vertices of the finite tree after subdivision.

The sequence of consecutive representative arcs can therefore cross any
edge of $H$ at most twice: it can enter and leave the consecutive block
of indices on its distal side only once each.  Summing edge lengths
gives
\begin{equation}\label{arm:chronological-tree-tour}
 \sum_{i=1}^{k-1}
 d_T\bigl(\pi(t_{v_i}),\pi(t_{v_{i+1}})\bigr)
 \le2\ell(H).
\end{equation}
This is the outer-face tour inequality, with its ordering verified
directly from interval minima.  It also constructs the relevant planar
order: root $H$ at its point nearest the root of $T$, order each child
subtree by its consecutive block of marked indices, and put repeated
marks of an internal vertex in the intervening corners.  The boundary
tour of this ordered tree encounters the marks chronologically and uses
each positive-length edge twice.  The preceding edge-count argument
does not require choosing any of those corners.

Equations~\eqref{arm:finite-contour-tree-distance},
\eqref{arm:connected-box-skeleton-length}, and
\eqref{arm:chronological-tree-tour} give
Equation~\eqref{arm:sorted-contact-total}.  Finally,
\[
 E_i=(h_{v_i}-\mu_i)+(h_{v_{i+1}}-\mu_i)
\]
is a sum of two nonnegative numbers.  It dominates their difference in
absolute value and each summand separately, proving
Equation~\eqref{arm:sorted-contact-walls}.  A singleton component has an
empty sum and needs no separate argument.
\end{proof}

\begin{corollary}[Discarding large comparison errors]
\label{arm:sorted-contact-discard}
Under the hypotheses of Lemma~\ref{arm:sorted-scalar-contacts}, for every
$\eta>0$ the number of consecutive comparisons with $E_i>\eta$, summed
over all components, is at most
\[
 \frac{6}{\eta}\sum_{v=1}^m\omega_v.
\]
All remaining comparisons have both representative height error and
wall error at most $\eta$.
\end{corollary}

\begin{proof}
Sum $\eta\mathbf 1_{\{E_i>\eta\}}\le E_i$ and apply
Equation~\eqref{arm:sorted-contact-total}.
\end{proof}

\begin{remark}\label{arm:sorted-contact-adaptedness}
The errors $E_i$ may depend on the whole contour.  The Lemma is a
deterministic implication, not a claim that these errors are measurable
with respect to the past of a sampling slot.  In a probabilistic use one
first keeps only comparisons whose error is at most a prescribed
deterministic tolerance.  The consequent wall may then be stated using
the earlier representative height and that tolerance.  Any conditioning
or enumeration needed to select the retained comparisons must still be
justified in the probability argument.
\end{remark}

\subsection{A noncrossing forest witnessing the contacts}

The preceding proof works for the original graph components, without
assuming that their witnesses form a noncrossing family.  For counting
purposes, the following separate construction gives a simultaneous
noncrossing witness, with possible mergers of components.

\begin{lemma}[Uncrossing and contraction]
\label{arm:noncrossing-contact-forest}
In the setting of Lemma~\ref{arm:sorted-scalar-contacts}, one may add
minimum chords whose endpoints are among the original chord endpoints
and then retain an outerplanar forest on the interval vertices with the
following properties.
\begin{enumerate}[label=\textup{(\roman*)}]
 \item The forest has one tree in each component of the enlarged chord
 graph.  Every original graph component is contained in one of these
 components.
 \item Every forest edge is witnessed by a genuine minimum chord.
 \item The vertices occur on the outer boundary in their time order.
 There are at most $32^m$ possible simple forests of this form on $m$
 prescribed ordered vertices.
\end{enumerate}
If the construction is applied to the endpoints belonging to a single
original connected component, it produces a noncrossing spanning tree
on exactly that component's vertices.  Forests obtained separately in
this way need not be noncrossing with one another.
\end{lemma}

\begin{proof}
On the finite set of original chord endpoints, put $s\sim t$ when
$f(s)=f(t)=M(s,t)$.  This is an equivalence relation, either directly
from the interval-minimum rule or from
Equation~\eqref{arm:finite-contour-tree-distance}.  For each class, add
the minimum chords between its consecutive times.  They connect all
times of that class, so they connect the endpoints of every original
chord.

These new chords can be drawn without crossings in the upper half-plane.
In fact, suppose two minimum chords have strictly interleaving endpoints
$a<b<c<d$, with pairs $(a,c)$ and $(b,d)$ at heights $h$ and $h'$.  Since
$b\in[a,c]$ and $c\in[b,d]$, the minimum conditions give respectively
$h'\ge h$ and $h\ge h'$.  Both heights agree, and all four times belong
to the same equivalence class.  Consecutive pairs of one class do not
interleave.  Chords from different classes therefore do not interleave
either.  Open semicircular arcs with nested or disjoint time intervals
as diameters do not intersect.  Shared
endpoints belong to the same class and cause no crossing.

Contract each interval $I_v$ in the boundary line to a single point.
This preserves an embedding with all vertices on the outer boundary.
One explicit way to see this is to choose a continuous proper
nondecreasing function $\varphi$ on the line that is constant precisely
on the nondegenerate $I_v$ and strictly increasing elsewhere.  For
$y>0$, the map
\[
 (x,y)\longmapsto
 \left(\frac{\varphi(x)+yx}{1+y},y\right)
\]
is a homeomorphism of the open upper half-plane, extends continuously
to the boundary, and there makes exactly the desired contractions.
Thus distinct chord interiors remain disjoint.  Discard loops and keep
one edge from each collection of parallel edges.  In every resulting
component choose a spanning tree, retaining its plane embedding.  The
earlier connectivity observation proves the first two assertions.
Using only endpoints from a single original connected component gives
the final assertion as well.

For completeness, the number of simple noncrossing forests on a fixed
ordered vertex set is exponential in its size.  Place the vertices on a
convex $m$-gon in that cyclic order.  Any noncrossing forest extends to a
triangulation after missing polygon sides are added and the remaining
faces are triangulated.  For $m\ge3$ there are fewer than $4^m$
triangulations, by the Catalan recursion, and each has fewer than $3m$
edges including the polygon boundary.  Its edge subsets number at most
$2^{3m}$.  Their product is at most $32^m$.  The cases $m=1,2$ satisfy
the same bound directly.
\end{proof}

\begin{remark}\label{arm:contact-mergers}
A simultaneous noncrossing forest can require merging original
components.  For example, two specified chords with interleaving
endpoints have a common contour height and allow additional true
chords between their endpoints, even if the specified graph initially
placed the two pairs in different components.  Such mergers are
explicit in Lemma~\ref{arm:noncrossing-contact-forest}; they must not be
silently treated as preservation of the original partition.  Adding
these true constraints cannot decrease the rank of the monochrome
relation.  Lemma~\ref{arm:sorted-scalar-contacts} itself needs no mergers.
\end{remark}
